\section{Discussion}
\label{sec:discussion}

\subsection{What the controls change}

RGBN has stronger release support, but absent construction provenance prevents
treating its Red index as verified. BGRN is a conservative sensitivity, not an
equally evidenced mapping. Index-versus-generic F1 directions survive this
audit; the index-modulated boundary direction does not.

Generic boundary weighting improves all 18 HR-GLDD comparisons, but no external
analysis confirms transfer. CAS is post hoc; LRD has weak validation, endpoint
instability, and trigger overlap. Sen12 fails four of five conditions, including
absolute performance. Its transformed-envelope integrity failure reflects
native rectangles that only abut, not duplicated samples. Transfer is
unsupported, not refuted.

This does not make NDVI universally unhelpful: prior work reports useful
spectral processing and NDVI guidance \citep{devara2024landslide,trigefnet2026}.
Our result is diagnostic, not a competitive model claim.

\subsection{Interpretation under incomplete observability}

Band-ratio cancellation is exact only for shared multiplicative irradiance.
Diffuse light, shadow, adjacency, reflectance, and atmosphere remain; an NDVI
gradient also responds to any vegetation edge. Without illumination or
background labels, the experiment cannot isolate scarps.

For incompletely observed remote-sensing systems, dataset-specific calibration
is essential \citep{varotsos2026information}. We likewise separate same-sensor
degradation from dataset shift and withhold event-level inference.

\subsection{Implications for benchmark practice}

Four recommendations follow: publish machine-readable channel names; bind
transfer and uncertainty claims to event/scene/spatial units; include generic
same-class controls (zero/raw-edge inputs and equal-mass plain boundary loss);
and propagate unresolved schemas into a machine-readable envelope.

The compact HR-GLDD model is not a transferable operational system; parameter
count alone says nothing about an emergency workflow.
